% !TEX root = ../main.tex
\section{Conclusion}\label{sec:CONCLUSION}
In this technical note, we have presented the near-threshold reconstruction algorithms for the \gls{dune} Phase~I \gls{fd} modules (\gls{hd} and \gls{vd}), focusing on the clustering, the vertex and the neutrino energy resolution of low-energy \gls{cc} interactions. The reconstruction is based on the \gls{tpc} cluster vertex reconstruction and \gls{tpc}--\gls{pds} flash matching, which allows us to extract the calibrated neutrino energy and the three-dimensional vertex position. The algorithm has been evaluated for \qtyrange{4}{30}{\mega\electronvolt} neutrinos, and optimised for the \qtyrange{4}{19}{\mega\electronvolt} solar neutrino range, in the presence of the full radiological and external background model (version 3), and demonstrates the first end-to-end low-energy workflow for these detectors, with more than \SI{90}{\percent} of the electrons from \gls{cc} signals above \SI{4}{\mega\electronvolt} reconstructed.

The vertex resolution was evaluated using a point--like reconstruction approach that combines the \gls{tpc} cluster vertex with the associated \gls{pds} flash. Characterised as the width of the core component of a double-Gaussian fit to the vertex error distribution, it is \SIrange{5}{7}{\milli\meter} for all coordinates and both geometries, with a broader component of about \SI{1.6}{\centi\meter}; the resolution grows with the neutrino energy following the electron track length. Including the tails, \SI{90}{\percent} (\SI{88}{\percent}) of the \gls{hd} (\gls{vd}) events are reconstructed within \SI{3}{\centi\meter} of the true $y$ and $z$ coordinates, and \SI{71}{\percent} (\SI{54}{\percent}) in $x$. The reconstruction of the $x$ coordinate additionally requires the \gls{tpc}--\gls{pds} flash matching, whose efficiency, defined as the fraction of events matched to a flash with a traceable signal contribution, averages \num{66}\% for \gls{hd} and \num{35}\% for \gls{vd}, reaching its minimum value at the furthest distance from the \gls{pds} plane. The efficiency of the $x$-coordinate vertex reconstruction itself, which also requires the reconstructed position to fall within the vertex-resolution tolerance, is lower: \num{54}\% (\num{27}\%) for \gls{hd} (\gls{vd}). Since the matching efficiency is lowest at the detector boundary, there are significant differences in the fiducialisation performance between the two geometries: the \gls{vd} longer drift distance and higher background rate result in a poorer ability to reject background events, while the \gls{hd} geometry achieves higher fiducialisation efficiency.

The neutrino energy was reconstructed using four different algorithms to sum the energy of the reconstructed \gls{tpc} clusters, evaluated in terms of energy resolution, fiducialisation efficiency, and the impact of flash matching on the energy reconstruction. With the \textit{SelectedEnergy} algorithm the energy resolution is most effective in the \qtyrange{6}{15}{\mega\electronvolt} range, reaching \SIrange{11}{14}{\percent} in \gls{hd} and \SIrange{16}{19}{\percent} in \gls{vd}, and \SI{14}{\percent} (\gls{hd}) and \SI{18}{\percent} (\gls{vd}) when averaged over \qtyrange{4}{20}{\mega\electronvolt}. The \gls{vd} energy resolution is slightly worse due to the higher background rate and the limited drift correction capabilities, which rely on the \gls{tpc}--\gls{pds} matching. The flash-matching step has a minor impact on the energy resolution compared to other factors, such as clustering inefficiencies and intrinsic energy losses. Overall, the reconstruction algorithms demonstrate good performance for low--energy neutrino interactions in \gls{dune}, with potential for further optimisation and improvement. Based on the amount of background activity in the \gls{vd} module, other reconstruction objectives such as channel tagging (out of the scope of this work) and particle identification are likely to be more challenging than in the \gls{hd} module. The case of a shielded \gls{vd} module, which has background rates comparable to the \gls{hd} module, is therefore highly encouraged and expected to yield similar low--energy reconstruction performance.

The reconstruction metrics presented here are based on the best estimates of the background and noise levels, and on the best models of the detector geometry, available at the time the studies were performed; they are expected to evolve as these models are revised and updated. The detector simulation and response modelled in this work are being validated in the context of the ProtoDUNE runs at CERN and will be fully proven once \gls{dune}'s Phase~I far detectors are operational.
